\documentclass{article}

\usepackage{luatexja} % This package is used for Japanese text in LuaLaTeX
\usepackage{amsmath}  % For better math handling
\usepackage{amssymb}  % For additional math symbols
\usepackage{graphicx}
\usepackage{listings}
\usepackage{amsthm}

\theoremstyle{definition}
\newtheorem{definition}{Definition}

\begin{document}
\subsection*{定義}
\begin{definition}
    マッチングgがg’にFirstOrderStochasticDominateされるとは、
    \begin{align*}
        \forall s \in S, c \in C,
        \sum_{c' \succ_s c} g_{s,c}(\succ) \leq \sum_{c' \succ_s c} g'_{s,c}(\succ)\\
        and \\
        \exists s \in S, c \in C,
        \sum_{c' \succ_s c} g_{s,c}(\succ) < \sum_{c' \succ_s c} g'_{s,c}(\succ)
    \end{align*}
\end{definition}
\begin{definition}
    メカニズムMがメカニズムM'にFirstOrderStochasticDominate(FOSD)されるとは、すべての選好プロファイルに対して、Mの出力マッチングがM'の出力マッチングにFirstOrderStochasticDominateされることである。すなわち、
    \begin{equation}
    \begin{aligned}
        \forall \succ \in P,\\
        &\forall s \in S, c \in C,
        &\sum_{c' \succ_s c} g_{s,c}(\succ) \leq \sum_{c' \succ_s c} g'_{s,c}(\succ)\\
        and \\
        &\exists s \in S, c \in C,
        &\sum_{c' \succ_s c} g_{s,c}(\succ) < \sum_{c' \succ_s c} g'_{s,c}(\succ)
    \end{aligned}
    \end{equation}
\end{definition}


\subsection*{効率性の証明}
ある正の実数ベクトル$(\lambda_s)_{s \in S} \in \mathbb{R}_{++}^{|S|}$をとり、
\begin{align*}
L^M((\lambda_s)_{s \in S}) = \sum_{\succ \in P}\sum_{s \in S} \lambda_s \sum_{c \in C} \sum_{c' \succ_s c} g^M_{s,c}(\succ)
\end{align*}
とした時、これを最大化するメカニズムMがFOSDされないことを示す。
対偶をとり、メカニズムMがM'にFOSDされるとき、
\begin{equation}
\begin{aligned}
    &\forall (\lambda_s)_{s \in S}\in \mathbb{R}_{++}^{|S|}\\ 
    &L^M((\lambda_s)_{s \in S}) < L^{M'}((\lambda_s)_{s \in S})
\end{aligned}
\end{equation}
が成り立つことを示す。
メカニズムMがM'にFOSDされるとき、(1)より、以下が成立。
\begin{align*}
    &\forall \succ \in P, \forall s \in S,\\
    &\sum_{c \in C} \sum_{c' \succ_s c} g^M_{s,c}(\succ) < \sum_{c \in C}\sum_{c' \succ_s c} g^{M'}_{s,c}(\succ)\\
    \rightarrow
    &\forall \succ \in P,\forall s \in S,\forall (\lambda_s)_{s \in S} \in \mathbb{R}_{++}\\
    &\lambda_s \sum_{c \in C} \sum_{c' \succ_s c} g^M_{s,c}(\succ) < \lambda_s \sum_{c \in C}\sum_{c' \succ_s c} g^{M'}_{s,c}(\succ)\\
    \rightarrow
    &\forall \succ \in P,\forall (\lambda_s)_{s \in S} \in \mathbb{R}_{++}^{|S|}\\
    &\sum_{s \in S}\lambda_s \sum_{c \in C} \sum_{c' \succ_s c} g^M_{s,c}(\succ) < \sum_{s \in S}\lambda_s \sum_{c \in C}\sum_{c' \succ_s c} g^{M'}_{s,c}(\succ)\\
    \rightarrow
    &\forall (\lambda_s)_{s \in S} \in \mathbb{R}_{++}^{|S|}\\
    &\sum_{\succ \in P}\sum_{s \in S}\lambda_s \sum_{c \in C} \sum_{c' \succ_s c} g^M_{s,c}(\succ) < \sum_{\succ \in P}\sum_{s \in S}\lambda_s \sum_{c \in C}\sum_{c' \succ_s c} g^{M'}_{s,c}(\succ)\\
\end{align*}
(2)が成り立つことが示された。


\end{document}